\chapter{Analisi dei rischi}
\label{cap:analisi-rischi}

% ===================================================================
%  RISCHI DI PROGETTO (organizzativi, tecnici, di pianificazione)
% ===================================================================

\section{Rischi di progetto}
\label{sec:rischi-progetto}

\subsection{Disponibilità e qualità dei dataset di valutazione}
\label{sub:rischio-dataset}
% Lo stato dell'arte è orientato all'anonimizzazione: scarseggiano benchmark pubblici
% pensati per la sola *detection* di PII. Rischio di non poter validare i risultati
% in modo riproducibile (cfr. problema evidenziato nelle linee guida).

\subsection{Complessità dell'integrazione multi-approccio}
\label{sub:rischio-integrazione}
% La pipeline prevede regex, NLP/NER (GLiNER, BERT) e potenzialmente AI generativa.
% Far convivere output eterogenei in un'unica pipeline modulare introduce rischi
% di incoerenza, duplicazione e complessità architetturale.

\subsection{Costo computazionale e tempi di elaborazione}
\label{sub:rischio-costo-computazionale}
% Le linee guida evidenziano che l'AI generativa su un intero filesystem
% richiederebbe "MESI di calcoli super computer-intensive". Necessario
% documentare il trade-off risultati/consumo energetico con paper accademici.

\subsection{Dipendenza da modelli e librerie di terze parti}
\label{sub:rischio-dipendenze}
% GLiNER, Hugging Face Transformers, librerie OCR, parser PDF/DOCX/XLSX.
% Rischio di breaking changes, discontinuità del supporto o licenze incompatibili.

\subsection{Carenza di competenze specialistiche}
\label{sub:rischio-competenze}
% Le linee guida sottolineano: "DEVO STUDIARE COSA SIGNIFICA TRAINING E COME
% FARLO IN MANIERA CORRETTA ED EFFICACE". Rischio legato alla curva di
% apprendimento su ML, fine-tuning e NLP nel contesto di un progetto di tesi.

\subsection{Scope creep e requisiti in evoluzione}
\label{sub:rischio-scope-creep}
% Molti requisiti sono ancora marcati "DA VALUTARE" (aggregazione PII,
% knowledge base, impieghi alternativi dell'AI). Rischio di continua
% espansione del perimetro durante lo sviluppo.

% ===================================================================
%  RISCHI DI PRODOTTO (qualità, affidabilità, sicurezza del sistema)
% ===================================================================

\section{Rischi di prodotto}
\label{sec:rischi-prodotto}

\subsection{Accuratezza insufficiente nella detection delle PII}
\label{sub:rischio-accuratezza}
% Il sistema non anonimizza, ma solo rileva: un'accuratezza bassa mina
% direttamente il valore del report finale per il DPO. Discutere metriche
% (precision, recall, F1) e soglie di accettabilità.

\subsection{Falsi positivi e impatto sull'usabilità}
\label{sub:rischio-falsi-positivi}
% Troppi falsi positivi rendono il report inutilizzabile: il DPO perde
% fiducia nel sistema e smette di consultarlo. Impatto diverso dai falsi
% negativi, da trattare separatamente.

\subsection{Falsi negativi e mancata rilevazione di PII critiche}
\label{sub:rischio-falsi-negativi}
% Un falso negativo su dati degni di particolare protezione (art. 5 nLPD,
% art. 9 GDPR) può avere conseguenze legali gravi. Rischio asimmetrico
% rispetto ai falsi positivi.

\subsection{Generalizzazione cross-domain (domain-silo problem)}
\label{sub:rischio-domain-silo}
% Come evidenziato dallo stato dell'arte (PIIBench, Jha et al.), i sistemi
% addestrati su un dominio falliscono su testi eterogenei. Il sistema deve
% funzionare su documenti aziendali di natura molto diversa.

\subsection{Supporto multilingue e multilocalizzazione}
\label{sub:rischio-multilingue}
% Il contesto nLPD + GDPR implica almeno IT, DE, FR, EN. Le regex sono
% lingua-dipendenti, i modelli NER hanno performance disomogenee tra lingue.
% Anche i formati delle PII cambiano (AVS vs CF vs SSN).

\subsection{Scalabilità su filesystem di grandi dimensioni}
\label{sub:rischio-scalabilita}
% Il requisito prevede la scansione ricorsiva di interi filesystem.
% Rischio di tempi proibitivi, consumo di memoria e gestione di formati
% eterogenei (PDF, DOCX, XLSX, immagini con OCR).

\subsection{Qualità e completezza dei ROPA in input}
\label{sub:rischio-qualita-ropa}
% Le linee guida notano che i ROPA aziendali possono essere "scritti in
% maniera sbrigativa" (es. "dati anagrafici" invece di campi specifici).
% Un ROPA povero degrada il confronto e il report di conformità.

\subsection{Conformità normativa e aggiornamento dei framework legali}
\label{sub:rischio-conformita-normativa}
% Il sistema deve allinearsi a GDPR e nLPD, ma anche potenzialmente a
% leggi americane e altri framework. Le normative evolvono: rischio che
% le categorie di PII implementate diventino incomplete o obsolete.

\subsection{Sicurezza e protezione dei dati durante l'analisi}
\label{sub:rischio-sicurezza-dati}
% Paradosso: per rilevare PII il sistema deve accedervi. Il registro
% interno, il report e le strutture intermedie contengono essi stessi
% dati personali. Rischio di data breach durante l'analisi stessa.

% ===================================================================
%  VALUTAZIONE DEI RISCHI
% ===================================================================

\section{Valutazione dei rischi}
\label{sec:valutazione-rischi}

\subsection{Criteri di classificazione}
\label{sub:criteri-classificazione}
% Definire le scale di probabilità e impatto utilizzate per classificare
% ciascun rischio (es. scala 1-5, qualitativa, o basata su standard ISO 31000).

\subsection{Matrice probabilità--impatto}
\label{sub:matrice-probabilita-impatto}
% Tabella/matrice che posiziona tutti i rischi identificati nelle sezioni
% precedenti in base alla loro probabilità di occorrenza e gravità dell'impatto.

% ===================================================================
%  STRATEGIE DI MITIGAZIONE
% ===================================================================

\section{Strategie di mitigazione}
\label{sec:mitigazione-rischi}
% Per ciascun rischio ad alta/media priorità, descrivere le contromisure
% adottate o pianificate (es. pipeline modulare, fallback regex, test
% incrementali, dizionari ad-hoc, sandboxing dei dati durante l'analisi).